%% ----------------------------------------------------------------------------
%%  §2.2  Proprietà di N, proprietà di Archimede, densità di Q
%% ----------------------------------------------------------------------------
\section{I numeri naturali, la proprietà di Archimede, la densità di \texorpdfstring{$\Q$}{Q}}
\label{sec:archimede}

Alla luce delle nozioni introdotte nel paragrafo precedente riportiamo di seguito alcune
proprietà dell'insieme dei numeri naturali e dell'insieme dei numeri razionali. Nell'insieme
$\R$, a partire dall'elemento $1$, si possono determinare gli elementi $2=1+1$, $3=2+1$ e
così via; tali elementi costituiscono l'insieme $\N=\{1,2,3,\dots\}$ dei numeri naturali
(\cref{sec:nzq}), i cui elementi sono ordinati secondo le relazioni $1<2<3<\cdots$. Il
principio di induzione (\cref{teo:induzione}) precisa il significato dei puntini: $\N$ è
il più piccolo sottoinsieme di $\R$ che contiene $1$ e che, insieme a ogni suo elemento
$n$, contiene anche $n+1$.

\subsection{Due proprietà caratteristiche di \texorpdfstring{$\N$}{N}}

Cominciamo con un lemma tecnico che rende rigorosa un'evidenza: tra due naturali
consecutivi non ci sono altri naturali.

\begin{lemma}[Struttura di $\N$]\label{lem:struttura-N}
\leavevmode
\begin{enumerate}
  \item Per ogni $n\in\N$ si ha $n\ge1$.
  \item Se $n\in\N$ e $n>1$, allora $n-1\in\N$.
  \item Se $n,m\in\N$ e $m>n$, allora $m\ge n+1$.
\end{enumerate}
\end{lemma}

\begin{schema}
\Ipotesi{$n,m\in\N$; \ il principio di induzione (\cref{teo:induzione}).}
\Tesi{(i) $n\ge1$; \ (ii) $n>1\implica n-1\in\N$; \ (iii) $m>n\implica m\ge n+1$.}
\end{schema}

\begin{strategia}
Tutte e tre le proprietà riguardano «tutti i numeri naturali», e l'unico strumento che
abbiamo per dimostrare qualcosa per tutti i naturali è il principio di induzione. Per
(i) e (ii) l'induzione è immediata. Per (iii) si fa induzione su $n$, tenendo $m$
generico: il caso $n=1$ si riduce a (i) e (ii); nel passo induttivo, se $m>n+1$ si
«toglie $1$» a $m$ (lecito per (ii)) e ci si riconduce all'ipotesi induttiva.
\end{strategia}

\begin{dimostrazione}
(i) La proprietà vale per $n=1$; se $n\ge1$ allora $n+1\ge1+1>1$ (si ricordi che $0<1$,
\cref{es:zero-minore-uno}). Per induzione vale per ogni $n$.

(ii) Sia $\mathcal{P}(n)$ la proposizione «$n=1$ oppure $n-1\in\N$». $\mathcal{P}(1)$ è
vera. Se vale $\mathcal{P}(n)$, allora $\mathcal{P}(n+1)$ è vera perché
$(n+1)-1=n\in\N$. Dunque $\mathcal{P}(n)$ vale per ogni $n$; se $n>1$, allora $n\neq1$, e
quindi $n-1\in\N$.

(iii) Sia $\mathcal{P}(n)$ la proposizione «per ogni $m\in\N$, se $m>n$ allora
$m\ge n+1$». Per $n=1$: se $m>1$, per (ii) $m-1\in\N$ e per (i) $m-1\ge1$, cioè $m\ge2$.
Supponiamo vera $\mathcal{P}(n)$ e sia $m\in\N$ con $m>n+1$. Allora $m>1$, quindi
$m-1\in\N$ e $m-1>n$; per l'ipotesi induttiva $m-1\ge n+1$, cioè $m\ge n+2$. Questo è
$\mathcal{P}(n+1)$.
\end{dimostrazione}

Con ragionamenti analoghi (per induzione) si prova che la somma e il prodotto di numeri
naturali sono numeri naturali; lo diamo per acquisito. L'insieme $\N$ gode di due
proprietà caratteristiche, che enunciamo e dimostriamo.

\begin{teorema}[Principio del buon ordinamento]\label{teo:buon-ordinamento}\index{principio del buon ordinamento}
Ogni sottoinsieme non vuoto di $\N$ ha minimo.
\end{teorema}

\begin{schema}
\Ipotesi{$A\subseteq\N$, \ $A\neq\emptyset$.}
\Tesi{Esiste $\min A$.}
\end{schema}

\begin{strategia}
Ragioniamo per assurdo: supponiamo che $A$ non abbia minimo e dimostriamo che allora $A$ è
vuoto. L'idea è di «risalire» i naturali uno alla volta con l'induzione, mostrando che
nessuno di essi può stare in $A$: se $1,2,\dots,n$ non stanno in $A$ e ci stesse $n+1$,
allora $n+1$ sarebbe il più piccolo elemento di $A$ (non ci sono naturali tra $n$ e
$n+1$, \cref{lem:struttura-N}(iii)) --- ma avevamo supposto che il minimo non ci fosse.
\end{strategia}

\begin{dimostrazione}
Supponiamo per assurdo che $A$ non abbia minimo, e sia $\mathcal{P}(n)$ la proposizione
«nessun naturale $k\le n$ appartiene ad $A$».

\emph{Base.} Se fosse $1\in A$, poiché $1\le a$ per ogni $a\in A$ (\cref{lem:struttura-N}(i)),
$1$ sarebbe il minimo di $A$: escluso. Dunque $1\notin A$; e poiché l'unico naturale
$k\le1$ è $k=1$ (per il \cref{lem:struttura-N}(i) e la proprietà antisimmetrica),
$\mathcal{P}(1)$ è vera.

\emph{Passo induttivo.} Supponiamo vera $\mathcal{P}(n)$. Se fosse $n+1\in A$, allora
ogni $a\in A$ verificherebbe $a>n$ (per $\mathcal{P}(n)$), quindi $a\ge n+1$ per il
\cref{lem:struttura-N}(iii): $n+1$ sarebbe il minimo di $A$, escluso. Quindi
$n+1\notin A$; e poiché i naturali $k\le n+1$ sono quelli $\le n$ e $n+1$ stesso
(ancora per il lemma), vale $\mathcal{P}(n+1)$.

Per il principio di induzione $\mathcal{P}(n)$ vale per ogni $n$; in particolare ogni
$n\in\N$ non appartiene ad $A$ (si applichi $\mathcal{P}(n)$ con $k=n$). Dunque
$A=\emptyset$, contro l'ipotesi.
\end{dimostrazione}

\begin{teorema}\label{teo:max-N}
Ogni sottoinsieme non vuoto di $\N$, limitato superiormente, ha massimo.
\end{teorema}

\begin{schema}
\Ipotesi{$A\subseteq\N$, \ $A\neq\emptyset$, \ $A$ limitato superiormente.}
\Tesi{Esiste $\max A$.}
\end{schema}

\begin{strategia}
Per il \cref{teo:esistenza-sup}, $A$ ha un estremo superiore $M$; dobbiamo provare che
$M\in A$. Per la caratterizzazione \eqref{eq:caratt-sup} con $\epsilon=1$, c'è un
elemento $n\in A$ «vicino» a $M$, precisamente $n>M-1$. Ora sfruttiamo il fatto che gli
elementi di $A$ sono naturali, dunque «distano almeno $1$» l'uno dall'altro: un elemento
di $A$ più grande di $n$ dovrebbe essere $\ge n+1>M$, impossibile. Quindi $n$ è il
massimo.
\end{strategia}

\begin{dimostrazione}
Sia $M=\sup A$, che esiste per il \cref{teo:esistenza-sup}. Per la
\cref{prop:caratt-sup} con $\epsilon=1$, esiste $n\in A$ con $n>M-1$. Se esistesse
$m\in A$ con $m>n$, per il \cref{lem:struttura-N}(iii) sarebbe $m\ge n+1>M$, contro il
fatto che $M$ è un maggiorante di $A$. Dunque $m\le n$ per ogni $m\in A$, e $n\in A$:
$n=\max A$.
\end{dimostrazione}

\begin{osservazione}
I due teoremi si estendono ai sottoinsiemi di $\Z$: \emph{ogni sottoinsieme non vuoto di
$\Z$ limitato inferiormente ha minimo, e ogni sottoinsieme non vuoto di $\Z$ limitato
superiormente ha massimo}. Basta «traslare» l'insieme dentro $\N$: se $A\subseteq\Z$ ha un
minorante $\ell$ e $k\in\N$ è tale che $k>-\ell$, l'insieme $\{a+k : a\in A\}$ è contenuto
in $\N$, e il suo minimo, diminuito di $k$, è il minimo di $A$ (l'esistenza di un tale $k$
è garantita dalla proprietà di Archimede, che dimostriamo subito e la cui dimostrazione
non usa questa osservazione, quindi non c'è circolarità). Per il massimo si applica il
risultato sul minimo all'insieme $-A=\{-a : a\in A\}$, che è limitato inferiormente se
$A$ è limitato superiormente.
\end{osservazione}

\subsection{La proprietà di Archimede}

Tenendo conto delle proprietà precedenti e del teorema di esistenza dell'estremo
superiore si può dimostrare il risultato seguente, che nel \cref{cap:numeri} abbiamo già
usato una volta (nella dimostrazione del \cref{teo:Q-non-completo}) aggirandolo.

\begin{teorema}[Proprietà di Archimede]\label{teo:archimede}\index{proprietà di Archimede}
Per ogni $x\in\R$ esiste $n\in\N$ tale che $n>x$.
\end{teorema}

\begin{schema}
\Ipotesi{$x\in\R$.}
\Tesi{$\esiste n\in\N : n>x$. \ Equivalentemente: $\N$ non è limitato superiormente,
$\sup\N=+\infty$.}
\end{schema}

\begin{strategia}
La tesi dice che $\N$ non ha maggioranti; ragioniamo per assurdo supponendo che ne abbia.
Allora, per il \cref{teo:esistenza-sup}, $\N$ avrebbe un \emph{minimo} maggiorante $M$.
Ma $\N$ è «invariante per traslazione di $1$»: se $n\in\N$ anche $n+1\in\N$. Quindi
$n+1\le M$ per ogni $n$, cioè $n\le M-1$ per ogni $n$: anche $M-1$ sarebbe un
maggiorante, e più piccolo di $M$. Assurdo. Si noti che l'ingrediente decisivo è la
completezza: senza l'esistenza del sup il ragionamento non parte.
\end{strategia}

\begin{dimostrazione}
Se la proprietà di Archimede fosse falsa, esisterebbe $x\in\R$ con $n\le x$ per ogni
$n\in\N$: l'insieme $\N$ sarebbe limitato superiormente e quindi, per il
\cref{teo:esistenza-sup}, dotato di estremo superiore $M\in\R$. In particolare, per ogni
numero naturale $n$ sarebbe $n\le M$. Poiché anche $n+1$ è un numero naturale,
risulterebbe $n+1\le M$, cioè $n\le M-1$, per ogni $n\in\N$; il che è assurdo perché
$M-1$ sarebbe un maggiorante di $\N$, contrariamente al fatto che $M$ è il più piccolo dei
maggioranti.
\end{dimostrazione}

In altre parole, la proprietà di Archimede afferma che l'insieme $\N$ dei numeri naturali
non è limitato superiormente. Ne seguono alcune conseguenze di uso continuo.

\begin{corollario}\label{cor:archimede}
\leavevmode
\begin{enumerate}
  \item Per ogni $\epsilon>0$ esiste $n\in\N$ tale che $\dfrac1n<\epsilon$; quindi
  $\inf\set{\dfrac1n : n\in\N}=0$.
  \item Se $a>0$ e $b\in\R$, esiste $n\in\N$ tale che $na>b$.
\end{enumerate}
\end{corollario}

\begin{schema}
\Ipotesi{(i) $\epsilon>0$; \ (ii) $a>0$, \ $b\in\R$.}
\Tesi{(i) $\esiste n\in\N : \frac1n<\epsilon$, \ e quindi $\inf\{1/n : n\in\N\}=0$; \
(ii) $\esiste n\in\N : na>b$.}
\end{schema}

\begin{strategia}
Entrambe le affermazioni sono la proprietà di Archimede applicata a un numero reale
scelto in modo opportuno: $1/\epsilon$ per la (i), $b/a$ per la (ii). Le disuguaglianze
richieste si ottengono poi moltiplicando o dividendo per numeri positivi, che non
cambiano il verso. Per l'estremo inferiore si usa la caratterizzazione
\eqref{eq:caratt-inf}.
\end{strategia}

\begin{dimostrazione}
(i) Per il \cref{teo:archimede} esiste $n>1/\epsilon$, cioè $1/n<\epsilon$. Poiché
$1/n>0$ per ogni $n$, lo zero è un minorante di $\{1/n\}$, e per quanto appena visto è
il massimo dei minoranti (\eqref{eq:caratt-inf}). (ii) Per il \cref{teo:archimede}
esiste $n>b/a$; moltiplicando per $a>0$ si ottiene $na>b$.
\end{dimostrazione}

\begin{idea}[Il significato fisico della proprietà di Archimede]
La forma (ii) del corollario è quella originale, geometrica: \emph{dati due segmenti, un
opportuno multiplo del più corto supera il più lungo}. Con un righello da $30$ cm si può
misurare (riportandolo un numero sufficiente di volte) la lunghezza di qualsiasi
corridoio; con un passo di integrazione $\Delta t$, per piccolo che sia, una simulazione
numerica raggiunge in un numero finito di passi qualsiasi istante finale. In una
«grandezza non archimedea» esisterebbero invece segmenti \emph{infinitesimi}, che
riportati un numero qualsiasi di volte non raggiungono mai un segmento dato: i numeri
reali non contengono quantità di questo tipo.
\end{idea}

\begin{notastorica}[Archimede, Eudosso e gli infinitesimi]
Il principio porta il nome di Archimede di Siracusa (287--212 a.C.), che lo enuncia come
postulato nel trattato \emph{Sulla sfera e il cilindro} e lo usa sistematicamente nel
\emph{metodo di esaustione}, con cui calcolò l'area del segmento di parabola e il volume
della sfera; nella prefazione alla \emph{Quadratura della parabola} lo stesso Archimede
osserva però che il lemma era già stato usato dai geometri precedenti, e la tradizione
lo attribuisce a Eudosso di Cnido (IV secolo a.C.); una sua forma compare nella definizione~4 del libro~V degli \emph{Elementi} di Euclide.
Per tutto il Seicento e il Settecento, Leibniz e i suoi successori calcolarono con gli
«infinitesimi», quantità più piccole di ogni numero positivo eppure non nulle: grandezze,
appunto, non archimedee. L'Analisi dell'Ottocento li bandì, sostituendoli con la nozione
di limite. Furono riabilitati in modo rigoroso solo nel 1961 da Abraham Robinson con
l'\emph{analisi non standard}, che lavora in un campo ordinato più grande di $\R$ nel quale
la proprietà di Archimede, per l'appunto, non vale.
\end{notastorica}

\subsection{La parte intera}

\begin{proposizione}[Parte intera]\label{prop:parte-intera}\index{parte intera}
Per ogni $x\in\R$ esiste uno e un solo $m\in\Z$ tale che
\begin{equation}\label{eq:parte-intera}
  m\le x<m+1.
\end{equation}
Il numero $m$ si chiama \emph{parte intera} di $x$ e si indica con $\lfloor x\rfloor$
(oppure con $[x]$).
\end{proposizione}

\begin{schema}
\Ipotesi{$x\in\R$.}
\Tesi{$\esiste!\,m\in\Z : m\le x<m+1$.}
\end{schema}

\begin{strategia}
Per l'esistenza, $m$ dev'essere «il più grande intero $\le x$»: lo otterremo come
massimo dell'insieme degli interi minori o uguali a $x$. Per non uscire da $\N$, dove
abbiamo i teoremi già pronti, trasliamo tutto di un naturale $k$ abbastanza grande
(fornito da Archimede) e usiamo invece il \emph{minimo} dei naturali che superano $x+k$.
Per l'unicità si osserva che due interi che soddisfano la \eqref{eq:parte-intera}
distano meno di $1$, e due interi distinti distano almeno $1$.
\end{strategia}

\begin{dimostrazione}
\emph{Esistenza.} Per il \cref{teo:archimede} esiste $k\in\N$ con $k>-x$, cioè $x+k>0$,
e (di nuovo per Archimede) l'insieme $T=\{j\in\N : j>x+k\}$ è non vuoto. Per il
\cref{teo:buon-ordinamento} $T$ ha minimo $j_{0}$. Allora $j_{0}>x+k$; inoltre
$j_{0}-1\le x+k$: se $j_{0}=1$ ciò segue da $x+k>0$, altrimenti $j_{0}-1\in\N$
(\cref{lem:struttura-N}(ii)) e $j_{0}-1\notin T$ perché è minore del minimo. Posto
$m=j_{0}-1-k\in\Z$, le due disuguaglianze $j_{0}-1\le x+k<j_{0}$ diventano $m\le x<m+1$.

\emph{Unicità.} Se anche $m'\in\Z$ verifica $m'\le x<m'+1$, sommando $m\le x<m+1$ con
$-m'-1<-x\le-m'$ si ottiene $-1<m-m'<1$. Ma un intero non nullo $h$ ha $\abs{h}\in\N$,
quindi $\abs{h}\ge1$ (\cref{lem:struttura-N}(i)); dunque $m-m'=0$.
\end{dimostrazione}

\begin{marginfigure}
\centering
\begin{tikzpicture}
\begin{axis}[width=5.4cm, height=4.6cm, xmin=-2.4, xmax=3.4, ymin=-2.6, ymax=3.4,
  axis lines=left, xtick={-2,-1,0,1,2,3}, ytick={-2,-1,0,1,2,3}]
  \foreach \k in {-2,...,2} {
    \addplot[curva, domain=\k:\k+0.97] {\k};
    \addplot[only marks, mark=*, mark size=1.2pt, blu] coordinates {(\k,\k)};
    \addplot[only marks, mark=o, mark size=1.2pt, blu, fill=white] coordinates {(\k+1,\k)};}
\end{axis}
\end{tikzpicture}
\caption{Il grafico della funzione parte intera $y=\lfloor x\rfloor$: una «scala» con
gradini chiusi a sinistra e aperti a destra.}
\label{fig:parte-intera}
\end{marginfigure}

Ad esempio $\lfloor 2{,}7\rfloor=2$, $\lfloor 3\rfloor=3$, $\lfloor-2{,}7\rfloor=-3$
(\emph{non} $-2$: $-3\le-2{,}7<-2$). In Python la parte intera è
\texttt{math.floor}, da non confondere con \texttt{int}, che tronca verso lo zero:
\texttt{int(-2.7)} restituisce \texttt{-2}.

\subsection{La densità dei razionali}

Dalla proprietà di Archimede si ricava che l'insieme $\Q$ dei numeri razionali è
\emph{denso} in $\R$.

\begin{definizione}[Insieme denso]\index{insieme!denso}
Un insieme $D\subseteq\R$ si dice \emph{denso} in $\R$ se per ogni coppia $a,b$ di numeri
reali con $a<b$ esiste un elemento di $D$ compreso fra $a$ e $b$.
\end{definizione}

\begin{teorema}[Densità dei numeri razionali]\label{teo:densita-Q}\index{densità!dei razionali}
L'insieme $\Q$ dei numeri razionali è denso in $\R$: per ogni $a,b\in\R$ con $a<b$
esiste $q\in\Q$ tale che $a<q<b$.
\end{teorema}

\begin{schema}
\Ipotesi{$a,b\in\R$, \ $a<b$.}
\Tesi{$\esiste m\in\Z,\ n\in\N : \ a<\dfrac{m}{n}<b$.}
\end{schema}

\begin{strategia}
Cerchiamo una frazione $m/n$ nell'intervallo $(a,b)$. L'idea è di scegliere prima il
\emph{denominatore} $n$ in modo che la «griglia» delle frazioni di denominatore $n$,
cioè i punti $\dots,-\frac1n,0,\frac1n,\frac2n,\dots$, abbia un passo $\frac1n$ più
piccolo della lunghezza $b-a$ dell'intervallo: allora almeno un punto della griglia deve
cadere dentro (\cref{fig:densita-Q}). Il denominatore lo fornisce Archimede ($n>1/(b-a)$).
Come numeratore prendiamo il primo intero che supera $na$, cioè $m=\lfloor na\rfloor+1$:
è il primo punto della griglia a destra di $a$, e non può «scavalcare» $b$ perché il passo
è troppo corto.
\end{strategia}

\begin{dimostrazione}
Per il \cref{teo:archimede} esiste $n\in\N$ con $n>\frac{1}{b-a}$; per cui risulterà
$nb-na>1$. Sia $m=\lfloor na\rfloor+1$ (\cref{prop:parte-intera}); allora
$m-1\le na<m$, e quindi
\[
  na<m=(m-1)+1\le na+1<na+(nb-na)=nb .
\]
Dalle disuguaglianze $na<m<nb$, dividendo per $n>0$, segue $a<\frac mn<b$.
\end{dimostrazione}

\begin{figure}[htb]
\centering
\begin{tikzpicture}[x=1.05cm, every node/.style={font=\footnotesize}]
  \draw[-Latex, grigio] (-0.3,0) -- (10.5,0);
  \foreach \k in {0,...,10} {\draw[grigio] (\k,0.09) -- (\k,-0.09);}
  \node[below=3pt, grigio] at (4,0) {$\frac{m-1}{n}$};
  \node[below=3pt, verde] at (5,0) {$\frac{m}{n}$};
  \node[below=3pt, grigio] at (6,0) {$\frac{m+1}{n}$};
  \fill[blu] (4.35,0) circle (2pt) node[above=3pt, blu] {$a$};
  \fill[blu] (5.7,0) circle (2pt) node[above=3pt, blu] {$b$};
  \fill[verde] (5,0) circle (2.3pt);
  \draw[grigio, decorate, decoration={brace, amplitude=3pt}] (7,0.25) -- (8,0.25) node[midway, above=3pt] {passo $\frac1n<b-a$};
  \draw[blu, thick] (4.35,-0.55) -- (5.7,-0.55);
  \node[blu, below] at (5.03,-0.6) {$b-a$};
\end{tikzpicture}
\caption{L'idea della dimostrazione del \cref{teo:densita-Q}: se il passo $1/n$ della
griglia delle frazioni di denominatore $n$ è più corto dell'intervallo $(a,b)$, almeno
un punto $m/n$ della griglia cade all'interno.}
\label{fig:densita-Q}
\end{figure}

\begin{corollario}[Densità degli irrazionali]\label{cor:densita-irrazionali}
Anche l'insieme $\R\setminus\Q$ dei numeri irrazionali è denso in $\R$.
\end{corollario}

\begin{schema}
\Ipotesi{$a,b\in\R$, \ $a<b$.}
\Tesi{$\esiste t\in\R\setminus\Q : a<t<b$.}
\end{schema}

\begin{strategia}
Non occorre ripetere la costruzione: basta «spostare» il problema. Traslato
l'intervallo $(a,b)$ di $-\sqrt2$, il teorema di densità vi trova un razionale $q$;
riportato indietro, $q+\sqrt2$ cade in $(a,b)$ ed è irrazionale. Al posto di $\sqrt2$
andrebbe bene qualsiasi irrazionale fissato.
\end{strategia}

\begin{dimostrazione}
Siano $a<b$. Applicando il \cref{teo:densita-Q} alla coppia $a-\sqrt2<b-\sqrt2$ si trova
$q\in\Q$ con $a-\sqrt2<q<b-\sqrt2$, cioè $a<q+\sqrt2<b$ (l'esistenza di $\sqrt2$ in
$\R$ è garantita dal \cref{teo:radice-n}); e $q+\sqrt2$ è irrazionale,
perché somma di un razionale e di un irrazionale (esercizio del \cref{sec:nzq}).
\end{dimostrazione}

\begin{osservazione}[Denso non vuol dire «pieno»]
Il \cref{teo:densita-Q} e il \cref{cor:densita-irrazionali} dicono che razionali e
irrazionali sono «mescolati» in modo inestricabile: in ogni intervallo, per piccolo che
sia, ci sono infiniti razionali e infiniti irrazionali. Eppure $\Q$ ha i «buchi» del
\cref{teo:Q-non-completo}: essere denso non significa essere completo. È una distinzione
sottile che si rivela profonda: Cantor dimostrò nel 1874 che i razionali si possono
mettere in fila e numerare (sono un'infinità \emph{numerabile}), mentre i reali no. In
un senso preciso, «quasi tutti» i numeri reali sono irrazionali.
\end{osservazione}

\begin{python}[La dimostrazione della densità, trasformata in un algoritmo]
# Il teorema di densità è costruttivo: la dimostrazione dice
# esattamente quali n e m scegliere. Traduciamola in codice, con
# l'aritmetica esatta delle frazioni.
from fractions import Fraction as F
import math

def razionale_tra(a, b):
    """Restituisce m/n con a < m/n < b, seguendo la dimostrazione."""
    assert a < b
    n = math.floor(1 / (b - a)) + 1     # Archimede: n > 1/(b - a)
    m = math.floor(n * a) + 1           # primo intero che supera n*a
    q = F(m, n)
    assert a < q < b                    # la tesi del teorema
    return q

print(razionale_tra(F(1, 3), F(1, 2)))         # tra 1/3 e 1/2
print(razionale_tra(math.sqrt(2), 1.4143))     # tra √2 e 1.4143
print(razionale_tra(math.pi, math.pi + 1e-6))  # intervallo minuscolo
\end{python}
\begin{risultato}
3/7
16363/11570
3141593/1000000
\end{risultato}

\begin{esercizio}
Dimostrare che tra due numeri reali distinti esistono \emph{infiniti} numeri razionali.
\begin{soluzione}
Siano $a<b$. Per il \cref{teo:densita-Q} esiste $q_{1}\in(a,b)$ razionale; applicandolo
di nuovo a $(a,q_{1})$ si trova $q_{2}\in(a,q_{1})$, poi $q_{3}\in(a,q_{2})$ e così via:
per induzione si costruisce una successione di razionali $q_{1}>q_{2}>q_{3}>\cdots$
tutti distinti e tutti in $(a,b)$. \emph{Scelta logica:} «infiniti» si dimostra
esibendo un procedimento che non si arresta mai; il teorema di densità è il «motore»
che a ogni passo produce un elemento nuovo.
\end{soluzione}
\end{esercizio}

\subsection{Approfondimento: i numeri reali in forma decimale}
\label{subsec:forma-decimale}

Fin dalla scuola siamo abituati a pensare un numero reale come un \emph{allineamento
decimale}, eventualmente illimitato: $\frac14=0{,}25$, $\frac13=0{,}333\ldots$,
$\sqrt2=1{,}41421356\ldots$. Negli assiomi del \cref{sec:assiomi}, però, di cifre non si
parla affatto, e i puntini non hanno per ora alcun significato preciso. Seguendo
\cite{giusti}, mostriamo che la rappresentazione decimale è una \emph{conseguenza} degli
assiomi: le cifre di un numero si ricavano dalla parte intera, e il numero si ricostruisce
dalle sue cifre con l'estremo superiore. Le serie numeriche, con le quali un allineamento
decimale si legge come una somma infinita, arriveranno solo nel \cref{cap:serie}: qui non
ne avremo bisogno. Ci limitiamo ai numeri $x\ge0$; per $x<0$ si rappresenta $-x$ e si
antepone il segno meno.

\begin{definizione}[Troncamenti e cifre decimali]\label{def:allineamento-decimale}\index{rappresentazione decimale}\index{cifre decimali}
Sia $x\ge0$. Per ogni intero $n\ge0$ il numero razionale
\begin{equation}\label{eq:troncamento}
  x_{n}=\frac{\lfloor10^{n}x\rfloor}{10^{n}}
\end{equation}
si chiama \emph{troncamento} di $x$ alla $n$-esima cifra decimale. Si pone poi
$c_{0}=\lfloor x\rfloor$ e, per ogni $k\in\N$,
\begin{equation}\label{eq:cifra-decimale}
  c_{k}=\lfloor10^{k}x\rfloor-10\,\lfloor10^{k-1}x\rfloor ;
\end{equation}
$c_{k}$ è la \emph{$k$-esima cifra decimale} di $x$, e si scrive
$x=c_{0}{,}c_{1}c_{2}c_{3}\ldots$ (\emph{allineamento decimale} di $x$), dove l'intero
$c_{0}$ si scrive nel modo usuale.
\end{definizione}

Ad esempio, per $x=\frac{27}{8}$ si trova $\lfloor x\rfloor=3$, $\lfloor10x\rfloor=33$,
$\lfloor100x\rfloor=337$, $\lfloor1000x\rfloor=3375$ e poi
$\lfloor10^{k}x\rfloor=3375\cdot10^{k-3}$; quindi $c_{1}=33-30=3$, $c_{2}=337-330=7$,
$c_{3}=3375-3370=5$ e $c_{k}=0$ per $k\ge4$: $\frac{27}{8}=3{,}375$. La formula
\eqref{eq:cifra-decimale} non fa che leggere la cifra delle unità dell'intero
$\lfloor10^{k}x\rfloor$: moltiplicare per $10^{k}$ sposta la virgola di $k$ posti verso
destra, e la parte intera butta via ciò che resta dopo la virgola.

\begin{teorema}[Esistenza dell'allineamento decimale]\label{teo:cifre-decimali}
Sia $x\ge0$, con troncamenti $x_{n}$ e cifre $c_{k}$ come nella
\cref{def:allineamento-decimale}. Allora:
\begin{enumerate}
  \item $c_{k}\in\{0,1,\dots,9\}$ per ogni $k\in\N$;
  \item per ogni $n\ge0$ (la somma è vuota se $n=0$)
  \begin{equation}\label{eq:troncamento-stima}
    x_{n}=c_{0}+\sum_{k=1}^{n}\frac{c_{k}}{10^{k}},
    \qquad x_{n}\le x<x_{n}+\frac{1}{10^{n}} ;
  \end{equation}
  \item $x=\sup\{x_{n} : n\ge0\}$;
  \item le cifre non sono \emph{definitivamente uguali a $9$}: non esiste $m\in\N$ tale
  che $c_{k}=9$ per ogni $k>m$.
\end{enumerate}
\end{teorema}

\begin{schema}
\Ipotesi{$x\ge0$; \ $x_{n}=\lfloor10^{n}x\rfloor/10^{n}$, \ $c_{0}=\lfloor x\rfloor$, \
$c_{k}=\lfloor10^{k}x\rfloor-10\lfloor10^{k-1}x\rfloor$.}
\Tesi{(i) $c_{k}\in\{0,\dots,9\}$; \ (ii) $x_{n}=c_{0}+\sum_{k=1}^{n}c_{k}10^{-k}$ e
$x_{n}\le x<x_{n}+10^{-n}$; \ (iii) $x=\sup_{n}x_{n}$; \ (iv) per nessun $m$ si ha
$c_{k}=9$ per ogni $k>m$.}
\end{schema}

\begin{strategia}
Tutto discende dalla doppia disuguaglianza che definisce la parte intera
(\cref{prop:parte-intera}), scritta per l'intero $m_{n}=\lfloor10^{n}x\rfloor$:
$m_{n}\le10^{n}x<m_{n}+1$. Dividendo per $10^{n}$ si ottiene la stima (ii):
\emph{il troncamento approssima $x$ per difetto con un errore minore di $10^{-n}$}.
Moltiplicando invece per $10$ la disuguaglianza scritta per $m_{k-1}$ si chiude $10^{k}x$
tra due interi che distano $10$, e questo costringe la cifra $c_{k}$ tra $0$ e $9$; la
formula di $x_{n}$ come somma delle cifre è una somma telescopica. Il punto (iii) è la
caratterizzazione dell'estremo superiore, con la proprietà di Archimede a garantire che
l'errore $10^{-n}$ diventi piccolo quanto si vuole. Il punto (iv) si prova per assurdo:
una coda di soli $9$ spingerebbe i troncamenti tanto vicino al numero $x_{m}+10^{-m}$ da
costringere $x$ a coincidere con esso, mentre la stima (ii) lo vieta.
\end{strategia}

\begin{dimostrazione}
Poniamo $m_{n}=\lfloor10^{n}x\rfloor\in\Z$, cosicché $x_{n}=m_{n}/10^{n}$ e, per la
\eqref{eq:parte-intera},
\begin{equation}\label{eq:mn-parte-intera}
  m_{n}\le10^{n}x<m_{n}+1 \qquad\text{per ogni } n\ge0 .
\end{equation}
Useremo un fatto elementare: \emph{se $j\in\Z$ e $j\le t$, allora $j\le\lfloor t\rfloor$}.
Infatti, se fosse $j>\lfloor t\rfloor$, l'intero $j-\lfloor t\rfloor$ sarebbe positivo,
quindi $\ge1$ (\cref{lem:struttura-N}(i)), e risulterebbe $j\ge\lfloor t\rfloor+1>t$.

(i) Moltiplicando per $10$ la \eqref{eq:mn-parte-intera} scritta con $n=k-1$ si ottiene
$10\,m_{k-1}\le10^{k}x<10\,m_{k-1}+10$. L'intero $10\,m_{k-1}$ non supera $10^{k}x$, e
quindi $10\,m_{k-1}\le m_{k}$; d'altra parte $m_{k}\le10^{k}x<10\,m_{k-1}+10$. Dunque
$c_{k}=m_{k}-10\,m_{k-1}$ è un intero con $0\le c_{k}<10$, cioè $c_{k}\in\{0,1,\dots,9\}$.

(ii) Dividendo la \eqref{eq:mn-parte-intera} per $10^{n}>0$ si ha
$x_{n}\le x<x_{n}+10^{-n}$. Inoltre $c_{k}10^{-k}=m_{k}10^{-k}-m_{k-1}10^{-(k-1)}
=x_{k}-x_{k-1}$, e la somma è telescopica:
\[
  c_{0}+\sum_{k=1}^{n}\frac{c_{k}}{10^{k}}=x_{0}+\sum_{k=1}^{n}(x_{k}-x_{k-1})=x_{n}.
\]

(iii) Per la (ii), $x$ è un maggiorante dell'insieme dei troncamenti. Osserviamo poi che
\begin{equation}\label{eq:dieci-meno-n}
  \perogni\epsilon>0\ \ \esiste n\in\N : \ \frac{1}{10^{n}}<\epsilon ,
\end{equation}
e che la stessa disuguaglianza vale allora per ogni indice più grande, perché le potenze
di $10$ crescono. Infatti per il \cref{cor:archimede}(i) esiste $n$ con $\frac1n<\epsilon$,
e per la disuguaglianza di Bernoulli (\cref{teo:bernoulli}) $10^{n}=(1+9)^{n}\ge1+9n>n$,
quindi $\frac{1}{10^{n}}<\frac1n<\epsilon$. Fissato $\epsilon>0$ e scelto $n$ come nella
\eqref{eq:dieci-meno-n}, la (ii) dà $x_{n}>x-10^{-n}>x-\epsilon$. Per la
\cref{prop:caratt-sup}, $x=\sup_{n}x_{n}$.

(iv) Supponiamo per assurdo che sia $c_{k}=9$ per ogni $k>m$, e poniamo $y=x_{m}+10^{-m}$.
Per $n>m$, usando la (ii) e l'identità $9\cdot10^{-k}=10^{-(k-1)}-10^{-k}$, che rende
telescopica anche questa somma,
\[
  x_{n}=x_{m}+\sum_{k=m+1}^{n}\frac{9}{10^{k}}=x_{m}+\frac{1}{10^{m}}-\frac{1}{10^{n}}
  =y-\frac{1}{10^{n}} .
\]
D'altra parte la (ii) con $n=m$ dà $x<y$, e con $n>m$ dà $x\ge x_{n}=y-10^{-n}$. Dunque
$0<y-x\le10^{-n}$ per ogni $n>m$, contro la \eqref{eq:dieci-meno-n} con $\epsilon=y-x$.
\end{dimostrazione}

\begin{figure}[htb]
\centering
\begin{tikzpicture}[x=0.8cm, every node/.style={font=\scriptsize}]
  % prima riga: [1, 2] diviso in decimi
  \draw[grigio] (0,0) -- (10,0);
  \foreach \k in {0,...,10} {\draw[grigio] (\k,0.08) -- (\k,-0.08);}
  \fill[rame!30] (4,-0.09) rectangle (5,0.09);
  \draw[rame, thick] (4,0) -- (5,0);
  \foreach \p/\lab in {0/1, 4/{1{,}4}, 5/{1{,}5}, 10/2} {\node[below=2pt] at (\p,0) {$\lab$};}
  \fill[blu] (4.142,0) circle (1.5pt) node[above=2pt] {$\sqrt2$};
  \node[anchor=west] at (10.3,0) {$x_{1}=1{,}4$, \ $c_{1}=4$};
  % seconda riga: [1,4; 1,5] diviso in centesimi
  \begin{scope}[yshift=-1.7cm]
  \draw[grigio] (0,0) -- (10,0);
  \foreach \k in {0,...,10} {\draw[grigio] (\k,0.08) -- (\k,-0.08);}
  \fill[rame!30] (1,-0.09) rectangle (2,0.09);
  \draw[rame, thick] (1,0) -- (2,0);
  \foreach \p/\lab in {0/{1{,}4}, 1/{1{,}41}, 2/{1{,}42}, 10/{1{,}5}} {\node[below=2pt] at (\p,0) {$\lab$};}
  \fill[blu] (1.421,0) circle (1.5pt) node[above=2pt] {$\sqrt2$};
  \node[anchor=west] at (10.3,0) {$x_{2}=1{,}41$, \ $c_{2}=1$};
  \end{scope}
  % terza riga: [1,41; 1,42] diviso in millesimi
  \begin{scope}[yshift=-3.4cm]
  \draw[grigio] (0,0) -- (10,0);
  \foreach \k in {0,...,10} {\draw[grigio] (\k,0.08) -- (\k,-0.08);}
  \fill[rame!30] (4,-0.09) rectangle (5,0.09);
  \draw[rame, thick] (4,0) -- (5,0);
  \foreach \p/\lab in {0/{1{,}41}, 10/{1{,}42}} {\node[below=2pt] at (\p,0) {$\lab$};}
  \node[below=2pt, anchor=north east] at (4.1,0) {$1{,}414$};
  \node[below=2pt, anchor=north west] at (4.9,0) {$1{,}415$};
  \fill[blu] (4.2136,0) circle (1.5pt) node[above=2pt] {$\sqrt2$};
  \node[anchor=west] at (10.3,0) {$x_{3}=1{,}414$, \ $c_{3}=4$};
  \end{scope}
  % la lente: l'intervallo evidenziato diventa la riga successiva
  \draw[grigio, dashed] (4,-0.5) -- (0,-1.45);
  \draw[grigio, dashed] (5,-0.5) -- (10,-1.45);
  \draw[grigio, dashed] (1,-2.2) -- (0,-3.15);
  \draw[grigio, dashed] (2,-2.2) -- (10,-3.15);
\end{tikzpicture}
\caption{Le cifre decimali come una lente d'ingrandimento. A ogni passo l'intervallo
$[x_{n},x_{n}+10^{-n})$ che contiene $x=\sqrt2$ (in color rame) viene diviso in dieci
parti uguali, e la cifra $c_{n+1}$ dice in quale delle dieci cade $x$. I troncamenti
$x_{n}$, estremi sinistri degli intervalli, crescono e hanno $x$ come estremo superiore
(\cref{teo:cifre-decimali}).}
\label{fig:decimali-lente}
\end{figure}

Il teorema precedente va dal numero alle cifre. Il cammino inverso, dalle cifre al numero,
è quello in cui interviene in modo essenziale la completezza di $\R$: assegnate
arbitrariamente le cifre, occorre \emph{trovare} un numero che le abbia.

\begin{teorema}[Ogni allineamento rappresenta un numero reale]\label{teo:allineamento-numero}
Siano $c_{0}\ge0$ un intero e $c_{1},c_{2},c_{3},\dots$ cifre di $\{0,1,\dots,9\}$, e
poniamo
\[
  s_{n}=c_{0}+\sum_{k=1}^{n}\frac{c_{k}}{10^{k}} \qquad (n\ge0).
\]
Allora l'insieme $\{s_{n} : n\ge0\}$ è limitato superiormente, e il suo estremo superiore
$s$ verifica $s_{n}\le s\le s_{n}+10^{-n}$ per ogni $n$. Inoltre:
\begin{enumerate}
  \item se le cifre $c_{k}$ non sono definitivamente uguali a $9$, l'allineamento decimale
  di $s$ è proprio $c_{0}{,}c_{1}c_{2}c_{3}\ldots$, ed $s$ è l'unico numero reale con
  questa proprietà;
  \item se invece $c_{k}=9$ per ogni $k>m$, allora $s=s_{m}+10^{-m}$, e le cifre di $s$
  sono nulle dalla $(m+1)$-esima in poi.
\end{enumerate}
\end{teorema}

\begin{schema}
\Ipotesi{$c_{0}\in\Z$, $c_{0}\ge0$; \ $c_{k}\in\{0,\dots,9\}$ per $k\in\N$; \
$s_{n}=c_{0}+\sum_{k=1}^{n}c_{k}10^{-k}$.}
\Tesi{Esiste $s=\sup_{n}s_{n}$, e $s_{n}\le s\le s_{n}+10^{-n}$; \ (i) se le cifre non
sono definitivamente $9$, le cifre di $s$ sono le $c_{k}$, ed $s$ è l'unico reale con
queste cifre; \ (ii) se $c_{k}=9$ per $k>m$, allora $s=s_{m}+10^{-m}$, con cifre nulle
dopo la $m$-esima.}
\end{schema}

\begin{strategia}
Le somme $s_{n}$ crescono, perché si aggiungono termini $\ge0$, ma non possono crescere
troppo: ciò che si aggiunge dopo l'indice $n$ è al più $\sum_{k>n}9\cdot10^{-k}$, una
somma telescopica minore di $10^{-n}$. Quindi $s_{n}+10^{-n}$ è un maggiorante, e il
candidato naturale per il numero cercato è $s=\sup_{n}s_{n}$: è la completezza di $\R$
(\cref{teo:esistenza-sup}) a garantire che \emph{esista}. Per il punto (i) si migliora la
stima in una disuguaglianza stretta, $s<s_{n}+10^{-n}$, sfruttando una cifra $c_{j}\le8$
con $j>n$; a quel punto $10^{n}s$ è stretto tra l'intero $10^{n}s_{n}$ e il suo
successivo, e l'unicità della parte intera dice che $s_{n}$ è il troncamento di $s$. Il
punto (ii) è un calcolo diretto, chiuso dalla caratterizzazione dell'estremo superiore.
\end{strategia}

\begin{dimostrazione}
Le $s_{n}$ crescono, perché $s_{n}-s_{n-1}=c_{n}10^{-n}\ge0$. Se $m>n$, poiché
$c_{k}\le9$ e $9\cdot10^{-k}=10^{-(k-1)}-10^{-k}$,
\begin{equation}\label{eq:coda-decimale}
  s_{m}-s_{n}=\sum_{k=n+1}^{m}\frac{c_{k}}{10^{k}}\le\sum_{k=n+1}^{m}\frac{9}{10^{k}}
  =\frac{1}{10^{n}}-\frac{1}{10^{m}}<\frac{1}{10^{n}} .
\end{equation}
Dunque $s_{m}<s_{n}+10^{-n}$ se $m>n$, e $s_{m}\le s_{n}$ se $m\le n$: il numero
$s_{n}+10^{-n}$ è un maggiorante di tutte le $s_{m}$. In particolare (con $n=0$)
l'insieme delle $s_{m}$ è limitato superiormente da $c_{0}+1$, e per il
\cref{teo:esistenza-sup} esiste $s=\sup_{m}s_{m}$; si ha $s\ge s_{n}$ perché $s$ è un
maggiorante, e $s\le s_{n}+10^{-n}$ perché $s$ è il minimo dei maggioranti.

(i) Fissato $n\ge0$, per ipotesi esiste $j>n$ con $c_{j}\le8$. Per $m\ge j$, ripetendo la
stima \eqref{eq:coda-decimale} ma usando $c_{j}\le9-1$ nel termine di indice $j$,
\[
  s_{m}-s_{n}\le\frac{1}{10^{n}}-\frac{1}{10^{m}}-\frac{1}{10^{j}}
  <\frac{1}{10^{n}}-\frac{1}{10^{j}} ;
\]
per $m<j$ si ha $s_{m}\le s_{j}$, e la stessa disuguaglianza vale ancora. Dunque
$s_{n}+10^{-n}-10^{-j}$ è un maggiorante delle $s_{m}$, e
\[
  s_{n}\le s\le s_{n}+\frac{1}{10^{n}}-\frac{1}{10^{j}}<s_{n}+\frac{1}{10^{n}} .
\]
Moltiplicando per $10^{n}$ si ottiene $10^{n}s_{n}\le10^{n}s<10^{n}s_{n}+1$, e
$10^{n}s_{n}=10^{n}c_{0}+\sum_{k=1}^{n}c_{k}10^{n-k}$ è un intero. Per l'unicità nella
\cref{prop:parte-intera}, $\lfloor10^{n}s\rfloor=10^{n}s_{n}$: il troncamento $n$-esimo
di $s$ è $s_{n}$. Allora $\lfloor s\rfloor=s_{0}=c_{0}$ e, per $k\in\N$, la $k$-esima
cifra di $s$ è $10^{k}s_{k}-10\cdot10^{k-1}s_{k-1}=10^{k}(s_{k}-s_{k-1})=c_{k}$. Se infine
$y\ge0$ ha le stesse cifre di $s$, per la \eqref{eq:troncamento-stima} i suoi troncamenti
sono $y_{n}=s_{n}$, e per il \cref{teo:cifre-decimali}(iii)
$y=\sup_{n}y_{n}=\sup_{n}s_{n}=s$.

(ii) Se $c_{k}=9$ per ogni $k>m$, la stima \eqref{eq:coda-decimale} è un'uguaglianza:
$s_{n}=s_{m}+10^{-m}-10^{-n}$ per $n>m$. Il numero $y=s_{m}+10^{-m}$ è quindi un
maggiorante delle $s_{n}$ (per $n\le m$ si ha $s_{n}\le s_{m}<y$) e, dato $\epsilon>0$,
scelto $n>m$ con $10^{-n}<\epsilon$ come nella \eqref{eq:dieci-meno-n}, si ha
$s_{n}>y-\epsilon$:
per la \cref{prop:caratt-sup}, $s=y$. Infine $10^{n}s=10^{n-m}(10^{m}s_{m}+1)$ è intero
per $n\ge m$, quindi $\lfloor10^{n}s\rfloor=10^{n}s$, e per $k>m$ la $k$-esima cifra di
$s$ è $10^{k}s-10\cdot10^{k-1}s=0$.
\end{dimostrazione}

\begin{osservazione}[Le code di $9$ e il legame con le serie]\label{oss:allineamenti-nove}
I due teoremi stabiliscono una corrispondenza biunivoca tra i numeri reali $x\ge0$ e gli
allineamenti decimali $c_{0}{,}c_{1}c_{2}\ldots$ le cui cifre non sono definitivamente
uguali a $9$. Gli allineamenti esclusi non sono «sbagliati»: per il
\cref{teo:allineamento-numero}(ii) rappresentano, tramite l'estremo superiore, numeri che
hanno anche un allineamento finito, ad esempio
\[
  0{,}4999\ldots=\sup\{0{,}4;\ 0{,}49;\ 0{,}499;\ \dots\}=0{,}5 ,
\]
e allo stesso modo $0{,}999\ldots=1$. Nel \cref{cap:serie} un allineamento sarà letto
come somma della serie $\sum_{k}c_{k}10^{-k}$: poiché le somme parziali crescono, la somma
è il loro estremo superiore (\cref{teo:monotone}), e i due punti di vista coincidono. Lì
vedremo anche che un allineamento periodico rappresenta un razionale
(\cref{prop:periodici}); il viceversa è l'\cref{es:decimali-periodici}. Infine, il numero
$10$ non ha nulla di speciale: in base $b\ge2$ le cifre stanno in $\{0,\dots,b-1\}$ e si
escludono le code di cifre $b-1$. I calcolatori usano la base $2$, in cui $\frac1{10}$ è
periodico: per questo \texttt{0.1} non è esattamente un decimo e
\texttt{0.1 + 0.2 == 0.3} è falso (\emph{Attenzione} alla fine del \cref{sec:assiomi}).
\end{osservazione}

\begin{python}[Le cifre di un numero, dalla definizione]
# Le cifre di x dalla definizione c_k = [b^k x] - b [b^(k-1) x],
# in aritmetica intera esatta: m(k) deve restituire [b^k x].
import math
from fractions import Fraction as F

def cifre(m, n, b=10):
    """Restituisce la stringa 'c_0,c_1c_2...c_n' nella base b."""
    c = [m(k) - b * m(k - 1) for k in range(1, n + 1)]
    return f"{m(0)}," + "".join(str(ck) for ck in c)

# [10^k sqrt(2)] è il più grande intero j con j^2 <= 2 * 10^(2k)
print("sqrt(2) =", cifre(lambda k: math.isqrt(2 * 100**k), 40))
# 1/7: le cifre si ripetono con periodo 6
print("1/7     =", cifre(lambda k: math.floor(F(10**k, 7)), 36))
# 1/10 in base 2: periodico, quindi 0.1 non è esatto in un float
print("1/10    =", cifre(lambda k: math.floor(F(2**k, 10)), 36, b=2))
print("valore esatto di 0.1:", F(0.1))
\end{python}
\begin{risultato}
sqrt(2) = 1,4142135623730950488016887242096980785696
1/7     = 0,142857142857142857142857142857142857
1/10    = 0,000110011001100110011001100110011001
valore esatto di 0.1: 3602879701896397/36028797018963968
\end{risultato}

Le cifre di $\sqrt2$ sono ottenute senza numeri in virgola mobile: per un intero
$j\ge0$, $j\le10^{k}\sqrt2$ equivale a $j^{2}\le2\cdot10^{2k}$, e \texttt{math.isqrt}
trova esattamente il più grande di tali $j$. È l'antica estrazione della radice «cifra
per cifra». L'ultima riga mostra il razionale (di denominatore $2^{55}$) che il
calcolatore memorizza al posto di $\frac1{10}$.

\begin{notastorica}[Stevin e l'«arte dei decimi»]
Le frazioni con denominatore $10$, $100$, $1000,\dots$ erano già note a matematici
cinesi e arabi, ma il loro uso sistematico in Europa si deve all'ingegnere fiammingo Simon
Stevin (1548--1620), che nel 1585 pubblicò a Leida un opuscolo di una trentina di pagine,
\emph{De Thiende} («Il decimo»), tradotto in francese lo stesso anno come \emph{La Disme}.
Lo dedicava ad astronomi, agrimensori, stazzatori di botti, zecchieri e mercanti: a
chiunque dovesse fare conti precisi con grandezze continue. Scriveva i decimali indicando
in un cerchietto l'esponente della potenza di un decimo, ad esempio $7{,}3486$ come
7\,\textcircled{\tiny 0}\,3\,\textcircled{\tiny 1}\,4\,\textcircled{\tiny 2}\,8\,%
\textcircled{\tiny 3}\,6\,\textcircled{\tiny 4}; la notazione, ripresa e semplificata da
Clavio e da Nepero, divenne quella attuale. Stevin prevedeva anche che monete, pesi e
misure sarebbero diventati decimali: accadde due secoli dopo, con il sistema metrico. Nella
sua \emph{Arithmétique}, dello stesso 1585, sosteneva infine che le radici e gli
irrazionali sono numeri a pieno titolo, e non quantità «assurde» o «sorde». Era, osserva
\cite{giusti}, la risposta che l'età dei mercanti dava all'incommensurabile scoperto dai
pitagorici: di $\sqrt2$ non si chiedeva che cosa fosse, ma come calcolarlo con tutte le
cifre necessarie. Ne resta traccia nella definizione scolastica di numero irrazionale come
«decimale illimitato non periodico», che prende per definizione ciò che abbiamo appena
dimostrato come teorema: le cifre esistono perché $\R$ è completo, non viceversa.
\end{notastorica}

\begin{esercizio}[I razionali hanno allineamento periodico]\label{es:decimali-periodici}
Sia $x=\frac pq$, con $p\ge0$ intero e $q\in\N$. Dimostrare che le cifre decimali di $x$
sono definitivamente periodiche: esistono $i\ge0$ e $\ell\in\N$ tali che
$c_{k+\ell}=c_{k}$ per ogni $k>i$.
\begin{soluzione}
Siano $m_{k}=\lfloor10^{k}x\rfloor$ e $r_{k}=10^{k}p-q\,m_{k}$: moltiplicando per $q$ la
\eqref{eq:mn-parte-intera} si trova $0\le r_{k}<q$, cioè $r_{k}$ è il resto della
divisione di $10^{k}p$ per $q$. Allora
\[
  10^{k+1}x=\frac{10\,(q\,m_{k}+r_{k})}{q}=10\,m_{k}+\frac{10\,r_{k}}{q},
\]
e poiché $\lfloor h+t\rfloor=h+\lfloor t\rfloor$ per ogni intero $h$ (per l'unicità della
parte intera), $m_{k+1}=10\,m_{k}+\lfloor10\,r_{k}/q\rfloor$. Dunque
\[
  c_{k+1}=\Bigl\lfloor\frac{10\,r_{k}}{q}\Bigr\rfloor, \qquad
  r_{k+1}=10\,r_{k}-q\,c_{k+1} :
\]
cifra e resto successivi dipendono \emph{soltanto} dal resto $r_{k}$: è la divisione in
colonna. I $q+1$ resti $r_{0},\dots,r_{q}$ stanno in $\{0,\dots,q-1\}$, che ha $q$
elementi, quindi due coincidono: $r_{i}=r_{i+\ell}$ con $\ell\ge1$. Da lì la divisione
ricomincia identica: per induzione $r_{k+\ell}=r_{k}$ per $k\ge i$, e quindi
$c_{k+\ell}=c_{k}$ per $k>i$. Per $\frac17$ i resti sono $1,3,2,6,4,5,1,\dots$ e il periodo
è $142857$. \emph{Scelta logica:} invece di una formula per le cifre, si cerca una
grandezza che ne determina il seguito (il resto) e che assume un numero finito di valori:
appena un valore si ripete, si ripete tutto («principio dei cassetti»).
\end{soluzione}
\end{esercizio}
